\section{不可微时次梯度为何仍能前进}
\label{sec:09-Convex-Optimization/06-Optimization-Algorithms/03}

给一个线性模型加上 \(\ell_1\) 正则，目的是把无关特征的系数压到零，好让下游同事看一张短的特征清单。损失里写上 \texttt{lam * w.abs().sum()}，交给 Adam 跑一万步。收敛了，验证集也不差，可打开权重一看：三千个系数里没有一个精确等于零，最小的那些停在 \(10^{-6}\) 到 \(10^{-8}\) 之间。想要的稀疏没有出现，得到的是一堆``几乎为零''。

问题出在绝对值的那个尖角上。\(\abs{w}\) 在 \(w=0\) 处没有导数，而自动微分必须返回点什么，于是框架按约定返回 \(0\)。这个返回值不算错，但迭代几乎永远踩不到 \(w=0\) 这个点：更新公式把 \(w\) 从正的一侧推过零点，落到负的一侧，下一步再推回来。真正想要的解就在尖角上，算法却只能在它两侧来回。

ReLU 也有同一个尖角，为什么那里没事？因为 ReLU 的折点通常不是我们要停留的地方，参数落在折点上是零测度事件，随便挑一个值当导数不影响大局。\(\ell_1\) 的折点恰恰是最优解所在之处——这是同一个数学现象，后果却完全不同。

\subsection{尖角处梯度不存在，支撑直线还在}

先看清尖角处到底缺了什么、还剩什么。可微点的梯度做两件事：它是最佳局部线性近似，同时是那条切线的斜率，而凸性保证这条切线是全函数的下界，见\hyperref[sec:09-Convex-Optimization/03-Convex-Functions/02]{``切平面为何是全局下界''}。在尖角处，第一件事塌了——左右两侧的极限不一致，最佳线性近似不存在。第二件事没塌：仍然有直线穿过这一点并且整条曲线都在它上方，只不过这样的直线不止一条。

把这些直线的斜率收集起来，就是次微分：

\[
\partial f(x)=\set{g\in\R^n:\ f(y)\ge f(x)+g^{\trans}(y-x)\ \text{对一切 }y}.
\]

定义里没有出现极限，也没有出现``近似''，只有一个对所有 \(y\) 成立的不等式。凸函数在其定义域内部处处有非空的次微分，可微点处这个集合退化成单点 \(\set{\nabla f(x)}\)。

\begin{center}
\begin{tikzpicture}[x=1.25cm,y=0.95cm,>=stealth]
  \draw[->,black!40] (-2.35,0) -- (2.5,0) node[right,font=\footnotesize] {\(x\)};
  \draw[->,black!40] (0,-1.4) -- (0,2.45);
  \foreach \s in {-1,-0.65,-0.3,0,0.3,0.65,1}
    {\draw[dashed,black!45] (-1.95,{-1.95*\s}) -- (1.95,{1.95*\s});}
  \draw[black!85,very thick] (-2.1,2.1) -- (0,0) -- (2.1,2.1);
  \fill[black] (0,0) circle (1.7pt);
  \node[font=\footnotesize,black!70] at (-1.5,2.3) {\(f(x)=\abs{x}\)};
  \node[below left,font=\footnotesize,black!55] at (-1.88,-1.9) {斜率 \(1\)};
  \node[below right,font=\footnotesize,black!55] at (1.88,-1.9) {斜率 \(-1\)};
  \begin{scope}[shift={(3.9,1.05)}]
    \draw[black!40] (-1.45,0) -- (1.45,0);
    \draw[black!85,very thick] (-1,0) -- (1,0);
    \fill[black] (-1,0) circle (1.6pt);
    \fill[black] (1,0) circle (1.6pt);
    \node[below,font=\footnotesize] at (-1,-0.06) {\(-1\)};
    \node[below,font=\footnotesize] at (1,-0.06) {\(1\)};
    \node[above,font=\footnotesize] at (0,0.1) {\(\partial f(0)=[-1,1]\)};
  \end{scope}
\end{tikzpicture}

\emph{原点处切线不唯一：任何斜率在 \(-1\) 与 \(1\) 之间的直线都从下方托住整条曲线，这一族斜率构成一个区间，而不是一个数。}
\end{center}

图上那束虚线全部穿过原点、全部落在图像下方，斜率填满 \([-1,1]\)。在 \(x>0\) 处这束线塌成一条，斜率只能取 \(1\)；越靠近原点，可选的斜率越多，到了原点就撑开成一整个区间。右边的线段把这个集合单独画了出来：不可微没有让信息消失，只是让答案从一个点变成一个集合。

集合形式的好处立刻显现。最优性条件从 \(\nabla f(x^\star)=0\) 变成

\[
0\in\partial f(x^\star),
\]

而对 \(f(x)=\abs{x}\)，\(0\) 确实落在 \([-1,1]\) 里，原点是最优点。把它用到 \(\ell_1\) 正则的目标 \(L(w)+\lambda\sum_j\abs{w_j}\) 上，第 \(j\) 个坐标取零当且仅当

\[
\abs{\frac{\partial L}{\partial w_j}(w)}\le\lambda .
\]

一个坐标只要在数据项上的推力不超过 \(\lambda\)，正则项就有足够的``摩擦''把它按在零上。稀疏不是四舍五入出来的，它写在这个不等式里，这也正是\hyperref[sec:09-Convex-Optimization/07-Applications/02]{``\(L_1\) 正则为何产生稀疏''}与\hyperref[sec:13-Machine-Learning/02-Linear-Models/02]{``正则化从 Ridge 到 Lasso''}讨论的几何在算法侧的对应。约束问题里把这个包含关系写全，就是\hyperref[sec:09-Convex-Optimization/05-Duality/03]{``KKT 把最优性写成一组方程''}的形式。

\subsection{负次梯度不保证下降，保证的是靠近}

有了次梯度，迭代格式可以照抄：任取 \(g_k\in\partial f(x_k)\)，令

\[
x_{k+1}=x_k-\alpha_kg_k .
\]

形状与梯度下降一样，性质却弱了一档。梯度下降之所以叫下降法，靠的是 \(-\nabla f(x)\) 是下降方向这一事实；而次梯度只承诺一个全局下界，不承诺沿 \(-g\) 走目标会变小。二维上很容易看出反例：在两个折面的交线上，某个次梯度可以垂直于交线的下行方向，沿它走一小段函数值反而上升。

那它凭什么还能前进？换一个被跟踪的量就清楚了。看到最优点的距离：

\[
\norm{x_{k+1}-x^\star}^2
=\norm{x_k-x^\star}^2-2\alpha_kg_k^{\trans}(x_k-x^\star)+\alpha_k^2\norm{g_k}^2 .
\]

次梯度不等式取 \(y=x^\star\) 给出 \(g_k^{\trans}(x_k-x^\star)\ge f(x_k)-f^\star\)。只要当前点还不是最优的，这个内积就严格为正，中间那一项就在把距离往下压；步长足够小时它盖过末项，距离严格减少。函数值可以上上下下，到最优点的距离却在稳步收紧——次梯度法沿着这条线索前进，而不是沿着目标值。

固定步长会在某个地方停住。回到 \(\abs{x}\)：\(x>0\) 时次梯度是 \(1\)，\(x<0\) 时是 \(-1\)，于是每一步的位移长度恒为 \(\alpha\)，一旦 \(\abs{x_k}<\alpha\)，迭代就越过零点跳到另一侧，此后永远在宽度约 \(\alpha\) 的区间里往返。这不是实现的缺陷，是固定步长的精度地板，形状与\hyperref[sec:09-Convex-Optimization/08-Nonconvex-and-Stochastic-Optimization/04]{``SGD 在期望中下降但噪声会留下平台''}里那个平台一致：那里地板由梯度噪声的方差决定，这里由次梯度不指向最优的那部分决定。想让误差继续掉，只能让 \(\alpha_k\) 递减。

\subsection{一阶信息变少，速度就掉下来}

把上面那个距离展开式对 \(k\) 求和，左侧望远镜相消，就得到全部的收敛信息。设初始距离不超过 \(R\)、次梯度范数不超过 \(G\)，把每个 \(f(x_k)-f^\star\) 松弛成其中的最小值，整理后是

\[
\min_{0\le k<N}\bigl(f(x_k)-f^\star\bigr)
\le
\frac{R^2+G^2\sum_{k<N}\alpha_k^2}{2\sum_{k<N}\alpha_k}.
\]

条件对一遍账：用到的只有次梯度不等式、次梯度有界、最优点存在，没有用到光滑性，也没有用到强凸。分子里的 \(G^2\sum\alpha_k^2\) 是每一步``走过头''的累计代价，分母里的 \(\sum\alpha_k\) 是累计行程，两者的比值就是能保证的精度。

这个式子马上给出步长该怎么选。若迭代预算 \(N\) 事先已知，取 \(\alpha_k=R/(G\sqrt N)\) 使分子分母平衡，界变成 \(RG/\sqrt N\)，也就是 \(O(1/\sqrt N)\)。若预算未知，就取满足 \(\sum_k\alpha_k=\infty\) 且 \(\sum_k\alpha_k^2<\infty\) 的序列，例如 \(\alpha_k\propto1/\sqrt{k+1}\)：前一个条件保证行程无限，能走到任意远；后一个条件保证过头的代价可积，不至于把已有的进展抵消掉。这两个条件是一对拉锯，减得太慢精度地板下不去，减得太快则可能在到达最优区域之前就失去移动能力。

若恰好知道最优值 \(f^\star\)，Polyak 步长

\[
\alpha_k=\frac{f(x_k)-f^\star}{\norm{g_k}^2}
\]

把展开式中的下降量做到最大，不需要任何手工调参。可行域简单时可以在更新后投影回去，\(x_{k+1}=\Pi_C(x_k-\alpha_kg_k)\)，因为投影不增加到可行点的距离，上面整套论证一字不改仍然成立，投影本身的几何见\hyperref[sec:03-Linear-Algebra/03-Orthogonality/02]{``投影与最近点：从一条直线走向一般子空间''}。

\(O(1/\sqrt N)\) 与光滑情形的 \(O(1/N)\) 差了一整档，而且这一档补不上：在只能访问次梯度的模型下，任何一阶方法在这个函数类上的复杂度下界就是 \(\Omega(1/\sqrt N)\)。慢不是算法笨，是我们把光滑性这个假设丢掉了，剩下的信息只够走这么快。一个附带的后果是，理论保证的对象是历史最优点或加权平均点，不是最后一个迭代点；实现时要么保留 best-so-far，要么按定理输出平均值，直接交出 \(x_N\) 会落在证明覆盖不到的地方。

\subsection{把不可微处当成黑箱，才慢成这样}

回到开头那个 \(\ell_1\) 的问题。次梯度法能给出正确的极限，却仍然交不出稀疏解：每一步都是一个连续的位移，落到坐标恰为零的点上是零概率事件，结果依旧是一堆 \(10^{-7}\)。慢和不稀疏是同一个根源——我们把 \(\abs{\cdot}\) 当成一个只能查询次梯度的黑箱，只用了它在当前点的一条支撑不等式，而它的全局形状明明是完全已知的。

hinge loss 的情形也一样，\(\max\set{0,1-yz}\) 的折点位置写在纸面上，见\hyperref[sec:09-Convex-Optimization/07-Applications/03]{``最大间隔与支持向量机''}。这类目标几乎总能拆成两半：一个光滑的数据项，加一个非光滑但结构极简单的正则项。对前者用梯度信息，对后者用它的解析结构，而不是只问一个次梯度——\hyperref[sec:09-Convex-Optimization/06-Optimization-Algorithms/04]{``近端与坐标方法利用可分结构''}从这里接手，并且会把 \(\ell_1\) 的速率重新拉回 \(O(1/k)\)。
